% conclusion.tex -- Enterprise Node Orchestrator Paper (RU-V5)

\section{Conclusion}

We presented the design and empirical evaluation of a pipelined, event-driven
state machine orchestrator for enterprise ZK validium nodes. The orchestrator
implements six states with 17 transitions, enabling concurrent transaction
ingestion during proof generation through three decoupled processing loops.

The key findings are:

\begin{enumerate}
\item \textbf{Orchestration overhead is negligible}: At 593~ms per
  8-transaction batch (0.66\% of the 90-second budget), the orchestration
  layer---including batch formation, SMT state updates, and witness
  generation---adds minimal overhead to the proof cycle. The overhead
  scales linearly at 73.5~ms/tx, entirely dominated by witness generation.

\item \textbf{Proving dominates latency}: ZK proof generation consumes
  64--85\% of end-to-end cycle time across all tested configurations.
  Performance optimization efforts should target prover backend selection
  and circuit constraint reduction, not orchestration engineering.

\item \textbf{Pipelining improves throughput modestly}: The 1.29$\times$
  speedup at batch 64 with rapidsnark reduces 10-batch processing by
  43~s. The speedup scales with the preparation-to-proving time ratio
  and would increase with faster provers or larger batches.

\item \textbf{Prover backend determines feasibility}: The snarkjs WASM
  backend fails the 90-second target at batch 64 (156.9~s), while
  rapidsnark meets it with 4.8$\times$ margin (18.9~s). Production
  deployment requires a native or GPU-accelerated prover for batch
  sizes above 16.
\end{enumerate}

The six formal invariants (liveness, safety, privacy, crash recovery,
state root continuity, single writer) provide a complete specification
for downstream TLA+ formalization. The component integration map---spanning
five prior research units covering SMT, circuits, L1 contracts, batch
aggregation, and data availability---establishes the implementation
blueprint for the Basis Network enterprise validium node.

Future work will replace mock components with production implementations
(Stage~2), validate crash recovery under fault injection (Stage~3), and
ablate individual components to quantify their marginal contribution to
system latency (Stage~4).
